\section{Setup and preliminaries}\label{sec:prelim}

\paragraph{Notation.}
We write $[N] := \{1,\dots,N\}$. For a function $u$ on $[0,1]^d$,
$\norm{u}_{\Ltwo}^2 := \int_{[0,1]^d} u(x)^2\,dx$ denotes its squared $\Ltwo$
norm, and for an integer $s\ge 1$ its squared Sobolev norm is
\begin{align}\label{eq:sobolev-norm}
\norm{u}_{\Sob}^2 \;:=\; \sum_{|\alpha|\le s} \norm{\partial^{\alpha}u}_{\Ltwo}^2,
\end{align}
the sum running over multi-indices $\alpha\in\N^d$ of order $|\alpha|\le s$,
with $\partial^{\alpha}$ the corresponding weak derivative. The Sobolev class
of interest is the unit ball
$\mathcal{W} := \{ f\in \Sob([0,1]^d) : \norm{f}_{\Sob}\le 1\}$.
For $\tau,\tau'\in\{0,1\}^{M_0}$, $\rhoH(\tau,\tau') := \sum_{j}\1\{\tau_j\ne\tau'_j\}$
is the Hamming distance and $\Nzero{\tau} := \sum_j \1\{\tau_j\ne 0\}$ the
Hamming weight. Throughout, $c, C, c_1, C_1, \dots$ denote universal positive
constants in the sense that they may depend only on the fixed quantities
$(s,d,g,\sigma^2)$ and on no other problem parameter (in particular not on the
sample size $n$ or the resolution $m$); their values may change from line to
line.

\begin{assumption}[Fixed-design Gaussian regression]\label{ass:setup}
We observe $D_n = \{(x_i,y_i)\}_{i=1}^n$ where the design points
$x_1,\dots,x_n$ are the fixed regular grid
$\{(\tfrac{i_1-1/2}{n^{1/d}},\dots,\tfrac{i_d-1/2}{n^{1/d}}) : i_\ell\in[n^{1/d}]\}$
in $[0,1]^d$ (so $n^{1/d}\in\N$), and
\begin{align}\label{eq:model}
y_i \;=\; f^*(x_i) + \xi_i, \qquad \xi_i \stackrel{\mathrm{iid}}{\sim} \mathcal{N}(0,\sigma^2),
\quad i\in[n],
\end{align}
for a known variance $\sigma^2>0$ and an unknown regression function
$f^*\in\mathcal{W}$. We write $P_{f}$ for the law of $D_n$ under regression
function $f$.
\end{assumption}

\begin{remark}\label{rem:setup}
The grid regularity is used only through \Cref{lem:kl}, where it lets us
compare the empirical sum $\frac1n\sum_i u(x_i)^2$ to $\norm{u}_{\Ltwo}^2$ for
the localized perturbations built in \Cref{def:bumps}; any design whose cells
each contain a comparable number of points would serve equally well. The
Gaussian-noise and known-$\sigma^2$ choices are what make the
pairwise Kullback--Leibler divergence in \Cref{lem:kl} a closed-form
mean-shift expression.
\end{remark}

\begin{fact}[Gaussian fixed-design KL]\label{fac:gauss-kl}
Under \Cref{ass:setup}, for any two regression functions $f,h$,
\begin{align}\label{eq:gauss-kl}
\KL\!\left(P_{f}\,\middle\|\,P_{h}\right)
\;=\; \frac{1}{2\sigma^2}\sum_{i=1}^{n}\bigl(f(x_i)-h(x_i)\bigr)^2.
\end{align}
\end{fact}
\begin{proof}
Under $P_f$ the observations are independent with $y_i\sim\mathcal{N}(f(x_i),\sigma^2)$,
and likewise under $P_h$; the designs coincide. Hence
\begin{align*}
\KL\!\left(P_{f}\,\middle\|\,P_{h}\right)
\;=\; & ~ \sum_{i=1}^{n}\KL\!\left(\mathcal{N}(f(x_i),\sigma^2)\,\middle\|\,\mathcal{N}(h(x_i),\sigma^2)\right) \\
\;=\; & ~ \sum_{i=1}^{n}\frac{\bigl(f(x_i)-h(x_i)\bigr)^2}{2\sigma^2},
\end{align*}
where the first equality is the tensorization of $\KL$ over the independent
coordinates and the second is the closed form for the $\KL$ between two
univariate Gaussians of common variance $\sigma^2$. This gives the assertion.
\end{proof}

\begin{definition}[Localized bump family]\label{def:bumps}
Fix once and for all a function $g\in C^\infty_c\bigl((0,1)^d\bigr)$ with
$g\not\equiv 0$, and set
$c_g := \norm{g}_{\Ltwo}^2 > 0$ and
$C_g := \sum_{|\alpha|\le s}\norm{\partial^{\alpha}g}_{\Ltwo}^2$.
For an integer resolution $m\ge 2$, partition $[0,1]^d$ into the $M_0 := m^d$
subcubes of side $1/m$, and let $x^{(j)}\in[0,1]^d$, $j\in[M_0]$, be the
"lower corners" so that $x\mapsto m\,(x-x^{(j)})$ maps the $j$-th subcube onto
$(0,1)^d$. For an amplitude $\omega>0$ and a codeword $\tau\in\{0,1\}^{M_0}$,
define
\begin{align}\label{eq:bump}
u_{\tau}(x) \;:=\; \sum_{j=1}^{M_0} \tau_j\,\frac{\omega}{m^{s}}\,
g\!\left(m\,(x-x^{(j)})\right), \qquad x\in[0,1]^d.
\end{align}
The summands have pairwise disjoint supports (one per subcube).
\end{definition}
